% Generated draft: requires mathematical and source review.
\documentclass[11pt]{article}
\usepackage[margin=25mm]{geometry}
\usepackage{fontspec}
\setmainfont{Latin Modern Roman}
\usepackage{amsmath,amssymb,mathtools,unicode-math}
\setmathfont{Latin Modern Math}
\usepackage{enumitem,needspace,xurl,hyperref}
\setlength{\parindent}{0pt}
\setlength{\parskip}{5pt}
\setlength{\emergencystretch}{3em}
\newcommand{\N}{\mathbb{N}}
\newcommand{\Z}{\mathbb{Z}}
\newcommand{\Q}{\mathbb{Q}}
\newcommand{\R}{\mathbb{R}}
\newcommand{\C}{\mathbb{C}}
\newcommand{\F}{\mathbb{F}}
\newcommand{\A}{\mathbb{A}}
\newcommand{\PP}{\mathbb{P}}
\newcommand{\E}{\mathbb{E}}
\newcommand{\eps}{\varepsilon}
\newcommand{\dd}{\,\mathrm{d}}
\begin{document}
\section*{UnsolvedMath ID: 30004751}
\textbf{Source problem code:} catalogue entry
\subsection{Definitions and mathematical statement}
Let \(G_{\Q}=\operatorname{Gal}(\overline{\Q}/\Q)\). For an elliptic curve \(E/\Q\) and a prime \(p\), let
\[
E[p]=\{P\in E(\overline{\Q}):pP=0\},
\]
viewed as a two-dimensional \(\F_p\)-vector space with its natural \(G_{\Q}\)-action. Two elliptic curves \(E,F/\Q\) are called \(p\)-congruent if there is a \(G_{\Q}\)-equivariant \(\F_p\)-linear isomorphism \(E[p]\cong F[p]\). No condition on the Weil-pairing multiplier is specified in the supplied record.

For \(d\in\Q^\times/(\Q^\times)^2\), write \(E^{(d)}\) for the quadratic twist of \(E\) by the square class \(d\). A simultaneous quadratic twist of a pair \((E,F)\) is \((E^{(d)},F^{(d)})\), using the same \(d\) for both curves. Write \([E,F]\) for the unordered pair of their \(\Q\)-isomorphism classes.

\textbf{Unresolved source qualifications.} The supplied record does not define ``nontrivial'' or identify the two known pairs by equations or other unique identifiers. Accordingly, let \(\mathcal T_p\) denote the intended collection of unordered \(\Q\)-isomorphism classes of nontrivial \(p\)-congruent pairs, with the meaning of nontrivial still to be supplied. Denote the two intended known pairs by \((E_1,F_1)\) and \((E_2,F_2)\); these are placeholders for specific pairs absent from the record.

\textbf{Classification question, conditional on resolving these qualifications.} Is it true that
\[
\mathcal T_{17}
=
\bigcup_{i=1}^{2}
\left\{[E_i^{(d)},F_i^{(d)}]:d\in\Q^\times/(\Q^\times)^2\right\},
\]
and that
\[
\mathcal T_p=\varnothing
\qquad\text{for every prime }p>17?
\]
The use of unordered pairs up to \(\Q\)-isomorphism makes explicit a comparison convention that is not stated in the supplied text. Until the two pairs and the intended nontriviality condition are specified, this is a conditional statement template rather than a fully identified classification conjecture.

The field labelled \texttt{original\_statement} instead concerns completeness of a list of exceptional images of elliptic curves over \(\Q\). It describes a purported equivalent condition involving unexpected rational points on specified modular curves associated with normalizers of nonsplit Cartan subgroups, or on one of two level-\(49\) modular curves of genus \(9\). Neither the exceptional-image list nor these curves are identified in the supplied record. No equivalence between that formulation and the \(p\)-congruence classification above is established by the supplied evidence.
\subsection{Short English statement}
The corrected source asks whether every nontrivial \(17\)-congruent pair of elliptic curves over \(\Q\) is a simultaneous quadratic twist of one of two known pairs, and whether no nontrivial \(p\)-congruent pairs exist for any prime \(p>17\). The two pairs and the meaning of nontrivial are missing. A conflicting original formulation concerns exceptional Galois images and modular curves, so the intended statement requires human source review.
\subsection{Sources}
\begin{itemize}
\item Supplied local registry record, permanent UnsolvedMath ID \(30004751\), title \emph{Pairs of \(p\)-Congruent Elliptic Curves}, exact source problem code \texttt{catalogue entry}, in the Oberwolfach Reports Open Problems collection. The record gives proposed year \(2021\) and names no proposer.
\item Supplied report citation: \emph{Explicit Methods in Number Theory (2022)}, Oberwolfach Reports, DOI \texttt{10.4171/owr/2021/34}, pp.~1805--1858. \url{https://doi.org/10.4171/owr/2021/34}. The citation's title year \(2022\), DOI year component \(2021\), and registry proposed year \(2021\) are retained as supplied.
\item Supplied literature source labelled ``Fisher'': \url{https://www.dpmms.cam.ac.uk/~taf1000/papers/congr17.pdf}. No paper title or fuller bibliographic details are supplied.
\item Supplied literature source labelled ``Oberwolfach report'': \url{https://ems.press/content/serial-article-files/46913}.
\item Literature assessment dated exactly \texttt{2026-08-22}: the supplied curation labels the problem open, says the classification for primes at least seventeen is not proved, and describes the conjectural list as the expected complete one in the primary work. This is supplied dated curation, not an independently verified literature assessment; the background explicitly requests independent verification of time-sensitive status and cited claims.
\item The registry labels the statement \texttt{corrected\_verified} and reports that it was checked against the official Oberwolfach report. These are supplied verification claims, not verification performed for this draft. The incompatible \texttt{original\_statement} remains an unresolved source issue.
\end{itemize}
\end{document}
